% Page 054

\seite{54}

hat bis Ende Oktober gewirtschaftet 20--23 Krieger verloren. Um in einigen\ob
Jahren so wurde auch in diesem Jahre der tapfern Krieger Weihnachten nicht\ob
vergessen, alle bekamen ein schönes Weihnachtspaket, was Ihnen besonders\ob
Freude gemacht, weil dieselben auf Anregung des Herrn Pfarrers besorgt wurden.

\pend
\abschnitt{Kriegsanleihe}
\pstart

An der vierten Kriegsanleihe beteiligten sich auch die Schulkinder.\ob
Es zeichneten 23 Kinder, zusammen 717 M.

\pend
\abschnitt{Schuljahr 1916-17}
\pstart

Mit dem 1.\ März wurden 4 Schüler entlassen, 2 Knaben, zwei\ob
Mädchen, diese widmen sich der Landwirtschaft, aufgenommen\ob
wurden 5 Kinder, 4 Mädchen, 1 Knabe. Die Schülerzahl be\ohb
trägt 33, davon sind 11 Knaben, 22 Mädchen.

\pend
\abschnitt{Ernte}
\pstart

In diesem Jahre war die Getreideernte eine bessere wie 1915, die Ernte\ob
selbst zog sich durch das schlechte Wetter sehr in die Länge, auch schon des\ohb
wegen, weil so wenige Arbeitskräfte hier sind. Kartoffeln gab es stellen\ohb
weise sehr wenig. Viehpreise waren den Sommer über hoch. Die Preise der Lebens\ohb
produkte sind noch fortwährend gestiegen. Eier, Butter 2,20 Mk – 2,40 Mk.

\pend
\abschnitt{Fortsetzung}
\pstart

Noch währt der furchtbare Krieg, aber das Volk bringt gerne weitere Opfer, die Land\ohb
bevölkerung bleibt daher nicht zurück. Nicht nur für die Soldaten im Felde, sondern\ob
auch der Armen in der Stadt wurde gedacht. Auf Anregung des Herrn Pfarrers wurden\ob
Sammlungen von Lebensprodukten, Fleisch, Fett, Eier, Kohl, Mehl im Orte gehalten\ob
und den Armen in Trier – später den Bedürftigen in heimatlichen Lazaretten über\ohb
wiesen. Zur Zeit der Ferien haben viele Kinder dieselben auf dem Lande zur\ob
Erholung verbracht. Die Schulkinder mußten auch mehr wie in vorigen Jahren bei\ob
der Feldarbeit helfen, darum hatten selbe den Sommer nur mit halben Tagen Schule und\ob
die Herbstferien wurden um 14 Tage verlängert. Es sind aus dem hiesigen\ob
Orte 9 gefallen.
\pend
